\subsubsection{Examples of Dual Spaces}\label{subsec:dual-space-examples}

\Needspace{4\baselineskip}
\begin{Example}
\label{ex:lp-conjugate-exponent-dual-functional}
Let $1<p<\infty$ and $p^{-1}+q^{-1}=1$. It follows from Hölder\textquotesingle s inequality that every $f\in L^p(\Omega)$ can be identified with a linear functional $l_f$ on $L^q(\Omega)$ by the correspondence
\begin{equation}\label{eq:lp-dual-functional-integral-pairing}
l_f(g):=\int_\Omega\overline{f(\bm{x})}g(\bm{x})\,\mathrm{d}\bm{x}.
\end{equation}
The complex conjugate is included to make the definition analogous to the inner product in Hilbert space. It is clear that $\|l_f\|\leq\|f\|_p$. Assume now that $f\ne0$ and let $f_n$ be a Cauchy sequence in $\widetilde L^p(\Omega)$ which converges to $f$ in $L^p$. Let $g_n=f_n|f_n|^{p-2}$. Then $g_n\in\widetilde L^q(\Omega)$ and $\|g_n\|_q^q=\|f_n\|_p^p$. Further, we find
\begin{equation}\label{eq:lp-dual-functional-approximating-norm-bound}
|l_f(g_n)-l_{f_n}(g_n)|\leq\|f-f_n\|_p\|g_n\|_q=\|f-f_n\|_p\|f_n\|_p^{p/q}
\end{equation}
and
\begin{equation}\label{eq:lp-dual-functional-approximating-pairing-value}
l_{f_n}(g_n)=\|f_n\|_p^p.
\end{equation}
It follows that $l_f(g_n)$ converges to $\|f\|_p^p$ and since $\|g_n\|_q$ converges to $\|f\|_p^{p/q}$, we conclude that actually $\|l_f\|=\|f\|_p$. Hence we have an isometric embedding from $L^p(\Omega)$ into the dual space of $L^q(\Omega)$. It can be shown that actually all functionals on $L^q(\Omega)$ can be represented in this way. We state this theorem without proof.
\end{Example}

\begin{Theorem}
\label{thm:lp-dual-space-representation}
Let $\Omega$ be an open, locally Jordan measurable set in $\mathbb{R}^m$, $1<p<\infty$ and $p^{-1}+q^{-1}=1$. Then the dual space of $L^p(\Omega)$ agrees with $L^q(\Omega)$ via \eqref{eq:lp-dual-functional-integral-pairing}.
\end{Theorem}

\Needspace{4\baselineskip}
\begin{Example}
\label{ex:bounded-continuous-lp-one-dual}
Let $f$ be a bounded continuous function on $\Omega$. Then we have
\begin{equation}\label{eq:bounded-continuous-lp-one-dual-estimate}
\left|\int_\Omega\overline{f(\bm{x})}g(\bm{x})\,\mathrm{d}\bm{x}\right|\leq\|g\|_1\sup_{\bm{x}\in\Omega}|f(\bm{x})|
\end{equation}
for every $g\in L^1(\Omega)$. Moreover, by taking $g$ to be non-negative with support near a point where $|f|$ is close to its supremum, we can see that
\begin{equation}\label{eq:bounded-continuous-infinity-norm}
\|f\|_\infty:=\sup_{\bm{x}\in\Omega}|f(\bm{x})|
\end{equation}
represents the norm of $f$ as a linear functional on $L^1(\Omega)$. However, not all linear functionals on $L^1(\Omega)$ can be represented in this way; for example we could take $f$\textquotesingle s which are discontinuous across some surface. We make the following definition.
\end{Example}

\begin{Definition}
\label{def:lp-infinity-dual-space}
Let $\Omega$ be an open, locally Jordan measurable set in $\mathbb{R}^m$. The dual space of $L^1(\Omega)$ is denoted by $L^\infty(\Omega)$.
\end{Definition}

Since test functions are dense in $L^1(\Omega)$, and convergence in the sense of test functions implies $L^1$-convergence, $L^\infty(\Omega)$ is a space of distributions. Moreover, since a uniform limit of continuous functions is also continuous, continuous functions cannot be dense in $L^\infty(\Omega)$; in fact, $L^\infty(\Omega)$ is not separable. In measure theory, $L^\infty$ is characterized as the space of bounded measurable functions.
